\begin{proposition}[Sharp triple packing for affine witnesses]
\label{prop:affine-witness-triple-packing}
Let $x_1,\ldots,x_N$ be distinct field elements and let $f,g$ be
received words. Suppose their common agreement with degree-$<2$
witnesses is at most $C$, where $2\le C<A\le N$.
Then the number of pairs $(z,h)$ with $\deg h<2$ and at least $A$
agreements with $f+zg$ is at most
\[
 \frac{N(N-1)(N-2)}{A(A-1)(A-C)}.
\]
In particular this bounds the number of qualifying challenges.
\end{proposition}
\begin{proof}
Write $h(X)=\alpha X+\beta$. Coordinate $i$ gives the affine plane
\[
 \alpha x_i+\beta-zg_i=f_i
\]
in parameter space $(\alpha,\beta,z)$. Any two distinct coordinates
give independent constraints and determine a line on which
$\alpha,\beta$ are affine functions of $z$. The coordinate planes
containing this entire line are exactly the common agreements of
the two corresponding affine witnesses for $f,g$, so there are at
most $C$ of them.

For any qualifying $(z,h)$, choose an ordered pair among $A$ of its
agreement coordinates. At least $A-C$ further agreement coordinates
have planes not containing the resulting line. These give at least
$A(A-1)(A-C)$ ordered triples of independent constraints. An independent
triple determines at most one $(z,h)$, and there are at most
$N(N-1)(N-2)$ ordered coordinate triples. Double counting proves the claim.
\end{proof}

\begin{corollary}[Asymptotic equality in the triple bound]
\label{cor:projective-linear-sharp-packing}
Let $p\ge5$ be prime, put
\[
 n=1+p+p^2+p^3+p^4,\qquad A=1+p+p^2,\qquad C=p+1,
\]
and use the domain $\mu_n\subseteq\mathbb F_{p^5}$.
Over $\mathbb F_{p^{15}}$, choose $\theta$ of degree three over
$\mathbb F_{p^5}$ and take
\[
 f(Y)=Y^{1+p+p^2+p^3}+\theta Y^{1+p+p^2}
          +\theta^2Y^{1+p},\qquad g(Y)=Y^{p+1}.
\]
Their common and individual agreements for affine witnesses equal $C$.
Exactly
\[
 M=\genfrac{[}{]}{0pt}{}{5}{2}_p=n(p^2+1)
\]
nonzero challenges have agreement greater than $C$; at every such
challenge the qualifying list is a singleton, with exact agreement $A$.
Moreover
\[
 M=\frac{n(n-1)(n-C)}{A(A-1)(A-C)}.
\]
Thus the universal upper bound of
Proposition~\ref{prop:affine-witness-triple-packing} is asymptotically
attained with constant one, and $M\sim n^{3/2}$.
\end{corollary}
\begin{proof}
Projectivize Proposition~\ref{prop:low-rate-large-characteristic-singletons}
by $Y=X^{p-1}$. Nonzero vectors of $\mathbb F_{p^5}$ have fibers of
size $p-1$, and its three-space agreement sets become projective planes
of size $A$. The witnesses become affine polynomials in $Y$.
Conversely an affine witness pulls back to a degree-$<p$ witness,
so the exact classification applies. The secondary agreements in that
proposition include zero; removing it and projectivizing gives exactly
$C$, not an additional high-agreement class.
The second source has degree $C$, and its projection in the first source
bounds both individual agreements by $C$. Projectivizing the simultaneous
witnesses on any two-space minus zero attains this bound.

There is also exact equality in the line-specific triple count.
Any two projective domain points span a projective line with $C$ points.
Their common interpolation pencil therefore lies in exactly $C$
coordinate planes. Hence there are exactly $n(n-1)(n-C)$ independent
ordered triples. Each such triple spans one projective plane, associated
with its unique three-space locator and challenge. Equivalently, the
displayed identity follows directly from
\[
 n-1=p(p+1)(p^2+1),\quad n-C=p^2A,\quad
 A-1=p(p+1),\quad A-C=p^2.
\]
The ratio to the universal upper bound is $(n-C)/(n-2)$, which tends
to one.
\end{proof}

The affine example is a sharp finite packing statement, rather than
an above-first-order asymptotic example: its agreement lies below the
rate-only curve $na_1(2/n)$ for all sufficiently large $p$.
The quadratic pullback in Theorem~\ref{thm:projective-quadratic-line}
retains the exponent $3/2$ while moving above that curve and below
exact Johnson at its advertised threshold.
